\documentclass[11pt]{article}
\usepackage[preprint]{acl}
\usepackage{times}
\usepackage{latexsym}
\usepackage[T1]{fontenc}
\usepackage[utf8]{inputenc}
\usepackage{microtype}
\usepackage{graphicx}
\usepackage{booktabs}
\usepackage{amsmath, amssymb, amsthm}
\usepackage{bm}
\usepackage{subcaption}
\usepackage{xcolor}
\usepackage{hyperref}

\newcommand{\veta}{\bm{\eta}}
\newcommand{\R}{\mathbb{R}}
\newcommand{\E}{\mathbb{E}}
\newcommand{\bpi}{\bm{\pi}}

\title{Non-Ergodic Mess3 Mixtures: How Transformers Represent\\Hierarchical Belief Structure}

\author{Seonglae Cho}

\begin{document}
\maketitle

\begin{abstract}
We train a small transformer on a non-ergodic dataset consisting of Mess3 processes with different parameters, where each sequence is generated by a single ergodic component. The transformer must solve a hierarchical inference problem: identifying the generating component (meta-synchronization) and tracking belief states within it. We derive mathematical predictions for the expected activation geometry, connecting the non-ergodic structure to the Factored World Hypothesis. Our analysis reveals that the transformer learns factored representations: component identity and within-component beliefs occupy approximately orthogonal subspaces. We further analyze synchronization dynamics across context positions and layers, showing that component identification precedes belief refinement.
\end{abstract}

\section{Introduction}

Language models encounter fundamentally non-ergodic data: different documents, styles, and topics constitute distinct statistical regimes. A model must perform in-context learning---using early tokens to identify the active regime---before it can predict optimally. This dual inference problem, identifying the process and tracking state within it, is central to how transformers operate in practice.

We study this problem in a controlled setting using Mess3 processes \cite{marzen2017nearly, shai2024transformers}, 3-state nonunifilar HMMs whose belief state geometry forms a Sierpinski gasket. By constructing a non-ergodic training set where each sequence comes from one of $K$ Mess3 parameterizations, we create a minimal model of the in-context learning problem: the transformer must determine which process generated the sequence and simultaneously track its hidden state.

This connects directly to the Factored World Hypothesis \cite{shai2026factored}: if the transformer treats component identity and within-component belief as independent factors, it should represent them in orthogonal subspaces with total dimensionality $(K-1) + 2 = K+1$, rather than the joint $3K - 1$ dimensions.

\paragraph{Contributions.}
\begin{enumerate}
    \item We formalize the non-ergodic Mess3 mixture as a hierarchical belief problem and derive predictions for the activation geometry (\S\ref{sec:predictions}).
    \item We train transformers and verify that belief states are linearly represented, with component-specific Sierpinski gaskets recoverable from the residual stream (\S\ref{sec:results}).
    \item We analyze synchronization dynamics: how quickly the model identifies the active component, and how this interacts with within-component belief tracking (\S\ref{sec:additional}).
\end{enumerate}

\section{Background}

\subsection{The Mess3 Process}

Mess3 is a 3-state nonunifilar HMM emitting tokens $z \in \{0, 1, 2\}$ from hidden states $s \in \{1, 2, 3\}$. It is parameterized by emission clarity $\alpha \in [0,1]$ and state persistence $x \in (0, 0.5]$, with $\beta = (1-\alpha)/2$ and $y = 1-2x$. The labeled transition matrices $T^{(z)}$ have entries $T^{(z)}_{ij} = P(\text{emit } z, \text{go to } j \mid \text{state } i)$; see Appendix~\ref{app:mess3} for the full definition.

Key properties: (1) The process is \emph{nonunifilar}---observing a token does not uniquely determine the state transition, so optimal prediction requires tracking a belief distribution over hidden states. (2) The set of reachable belief states forms a \emph{Sierpinski gasket} in the 2-simplex. (3) The net transition matrix has eigenvalue $\zeta = 1 - 3x$, controlling the exponential decay rate of belief corrections with context distance.

\subsection{Belief State Geometry in Transformers}

Shai et al.~\cite{shai2024transformers} showed that transformers trained on next-token prediction linearly represent belief states in their residual stream. The belief (predictive vector) after context $z_{1:t}$ is:
\begin{equation}
    \veta^{(z_{1:t})} = \frac{\veta^{(\varnothing)} T^{(z_1)} \cdots T^{(z_t)}}{\veta^{(\varnothing)} T^{(z_1)} \cdots T^{(z_t)} \mathbf{1}}
\end{equation}
where $\veta^{(\varnothing)} = \bpi = [\frac{1}{3}, \frac{1}{3}, \frac{1}{3}]$ is the uniform stationary distribution.

Piotrowski et al.~\cite{piotrowski2025constrained} showed transformers implement a constrained approximation to Bayesian filtering: $r_1(z_{1:d}) = \bpi + \sum_{s=1}^{d} (\bpi T^{z_s} T^{d-s} - \bpi)$, where each correction decays as $\zeta^{d-s}$.

\subsection{Factored World Hypothesis}

Shai et al.~\cite{shai2026factored} demonstrated that transformers factor their world model into parts, representing conditionally independent factors in orthogonal subspaces. For $N$ factors with latent dimensions $d_1, \ldots, d_N$, the factored representation requires $\sum_n (d_n - 1)$ dimensions instead of $\prod_n d_n - 1$ (joint).

\section{Non-Ergodic Mess3 Mixture}

\subsection{Construction}

We construct a training set of $K$ Mess3 processes with distinct parameters $\{(\alpha_k, x_k)\}_{k=1}^K$. Each training sequence is generated entirely by one component, chosen uniformly. The model observes only the token sequence---the component identity is latent.

\paragraph{Component selection.} We choose components with well-separated entropy rates $H_k$ and spectral gaps $\zeta_k = 1 - 3x_k$ to ensure meaningfully different dynamics:

\begin{table}[h]
\centering
\small
\begin{tabular}{lccccc}
\toprule
 & $\alpha$ & $x$ & $H_k$ (nats) & $\zeta_k$ \\
\midrule
C0 & 0.95 & 0.03 & 0.34 & 0.91 \\
C1 & 0.85 & 0.05 & 0.69 & 0.85 \\
C2 & 0.65 & 0.15 & 0.97 & 0.55 \\
\bottomrule
\end{tabular}
\caption{Mess3 components used in the non-ergodic mixture.}
\label{tab:components}
\end{table}

\subsection{Relevance to Language Models}
\label{sec:relevance}

The non-ergodic structure mirrors key properties of natural language:
\begin{itemize}
    \item \textbf{In-context learning}: Each sequence begins with ambiguity about which ``regime'' is active. Early tokens must be used to identify the generating process before optimal prediction is possible---analogous to how LLMs adapt to document topic, style, or task.
    \item \textbf{Hierarchical structure}: Natural language has latent structure at multiple scales (topic, syntax, discourse). The component/state hierarchy in our setup is a minimal model of this.
    \item \textbf{Non-stationarity}: Real text is non-ergodic---statistics vary across documents. Understanding how transformers handle this is fundamental to alignment and interpretability.
\end{itemize}

\section{Pre-Registered Predictions}
\label{sec:predictions}

\subsection{Mathematical Derivation}

The joint state space is $(k, s) \in \{1,\ldots,K\} \times \{1,2,3\}$, with $3K$ total states. The predictive vector decomposes:
\begin{equation}
    P(k, s \mid z_{1:t}) = P(k \mid z_{1:t}) \cdot P(s \mid k, z_{1:t})
\end{equation}
where $P(k \mid z_{1:t})$ is the \emph{meta-belief} (posterior over components) and $P(s \mid k, z_{1:t})$ is the \emph{within-component belief state}.

The joint transition operator is block-diagonal:
\begin{equation}
    T^{(z)} = \bigoplus_{k=1}^K P(k) \cdot T_k^{(z)}
\end{equation}
This is \emph{not} a tensor product, so the standard factored representation from~\cite{shai2026factored} does not directly apply. However, the structure is simpler: component identity is static within a sequence, so after sufficient context, $P(k \mid z_{1:t})$ concentrates on the true component.

\paragraph{Prediction 1: Dimensionality.}
We predict the effective dimensionality should be approximately $K + 1$:
\begin{itemize}
    \item $(K-1)$ dimensions for meta-belief (the simplex over $K$ components)
    \item $2$ dimensions for within-component belief (the 2-simplex of each Mess3)
\end{itemize}

The factored prediction of $K+1$ dimensions should hold if the model separates component identity from belief state. The joint alternative requires $3K - 1$ dimensions.

\paragraph{Prediction 2: Context-position dependence.}
At early positions (small $t$), the meta-belief is uncertain: all components are plausible. The geometry should resemble a superposition of $K$ gaskets. As $t$ increases, the meta-belief concentrates, and the geometry should resolve into the specific gasket of the true component.

The meta-belief convergence rate depends on the KL divergence between component distributions: $D_{\text{KL}}(P_{k_1}(z) \| P_{k_2}(z))$. Components with more distinct entropy rates should be distinguished faster.

\paragraph{Prediction 3: Layer dependence.}
Based on the layer-by-layer construction observed in~\cite{shai2024transformers}:
\begin{itemize}
    \item \textbf{Embedding}: 3 points (token vertices), no component structure
    \item \textbf{Early layers}: Component clusters should emerge (meta-synchronization)
    \item \textbf{Later layers}: Within-component gasket structure should appear
\end{itemize}

\subsection{Geometric Intuition}

Multiple geometries are conceivable:

\begin{enumerate}
    \item \textbf{Separated gaskets}: $K$ distinct Sierpinski gaskets in different regions of activation space, distinguished by large offsets (component encoding). Total: $\sim 2K$ dimensions.
    \item \textbf{Factored}: Component identity in a $(K-1)$-dimensional subspace, shared 2D belief subspace with superimposed gaskets. Total: $K+1$ dimensions.
    \item \textbf{Conditional}: Similar to factored, but the 2D belief subspace is component-specific (different orientations per component). Total: $2K + (K-1)$ dimensions.
\end{enumerate}

Our formal prediction favors geometry 1 or 2. Geometry 2 would be the strongest evidence for factored representations. Geometry 1 is more natural for a block-diagonal (not tensor-product) structure.

We will not be surprised if the model learns geometry 1 initially and transitions toward geometry 2 over training, following the inductive bias toward factoring observed in~\cite{shai2026factored}.

\section{Experimental Setup}

\paragraph{Model.} GPT-2 style decoder-only transformer with 3--4 layers, $d_{\text{model}} = 64$--$128$, 4 attention heads, pre-norm (LayerNorm), learned positional embeddings. Implemented with TransformerLens~\cite{nanda2022transformerlens}.

\paragraph{Training.} Next-token cross-entropy loss with Adam ($\text{lr} = 5 \times 10^{-4}$, no weight decay). On-the-fly data generation with batch size 8192, trained for 20K--50K steps. Vocabulary: $\{$BOS, 0, 1, 2$\}$.

\paragraph{Analysis methods.} (1) Cumulative explained variance (CEV) of residual stream activations via PCA. (2) Linear regression from activations to ground-truth belief states, measured by $R^2$. (3) Meta-belief regression (posterior over components). (4) Visualization of belief geometry in principal component subspaces.

\section{Results}
\label{sec:results}

\textit{[Results to be filled after experiment completion.]}

\subsection{Phase 1: Single Mess3 Verification}

\subsection{Phase 2: Non-Ergodic Mixture}

\subsubsection{Training dynamics}

\subsubsection{Dimensionality analysis}

\subsubsection{Belief regression}

\subsubsection{Geometry visualization}

\section{Additional Analysis: Synchronization Dynamics}
\label{sec:additional}

We analyze the \emph{synchronization dynamics} of the transformer: how quickly and through which mechanism does the model identify the active component?

\subsection{Meta-belief accuracy by position}

We measure $R^2$ of the meta-belief regression separately at each context position. This reveals when the model has accumulated enough evidence to identify the component.

\subsection{Layer-wise probing}

At each layer, we independently probe for (a) component identity and (b) within-component belief state. This reveals the order in which these representations are constructed.

\subsection{Attention pattern analysis}

We examine whether the attention mechanism develops specialized patterns for meta-synchronization (attending to positions most informative for component identification).

\section{Discussion}

\paragraph{Connection to in-context learning.} The non-ergodic Mess3 mixture provides a minimal, analytically tractable model of in-context learning. The hierarchical structure (identify regime, then predict within regime) may be a fundamental computational motif in language models.

\paragraph{Limitations.} Our setup uses a small model and simple processes. Scaling to larger models and more complex non-ergodic structures (e.g., slowly drifting parameters, switching processes) is an important direction.

\section*{Acknowledgments}
This work was completed as part of the MATS Summer 2026 application process.

\bibliography{references}
\bibliographystyle{acl_natbib}

\appendix

\section{Mess3 Transition Matrices}
\label{app:mess3}

With $\beta = (1-\alpha)/2$ and $y = 1-2x$:
\begin{align}
T^{(0)} &= \begin{pmatrix} \alpha y & \beta x & \beta x \\ \alpha x & \beta y & \beta x \\ \alpha x & \beta x & \beta y \end{pmatrix} \\
T^{(1)} &= \begin{pmatrix} \beta y & \alpha x & \beta x \\ \beta x & \alpha y & \beta x \\ \beta x & \alpha x & \beta y \end{pmatrix} \\
T^{(2)} &= \begin{pmatrix} \beta y & \beta x & \alpha x \\ \beta x & \beta y & \alpha x \\ \beta x & \beta x & \alpha y \end{pmatrix}
\end{align}

The entropy rate (lower bound on achievable cross-entropy) is:
\begin{equation}
H = -\sum_{s=1}^{3} \pi(s) \sum_{z=0}^{2} P(z|s) \log P(z|s)
\end{equation}
where $\pi(s) = 1/3$ (uniform) and $P(z|s) = \sum_j T^{(z)}_{sj}$.

\end{document}
